% ============================================================================
% Fichier  : partie1/P01_C01_hacking_ethique.tex
% Titre    : Le Hacking Éthique : Définitions, Métiers, Types de Tests
% Auteur   : MINKA MI NGUIDJOÏ Thierry Emmanuel
% Version  : 1.0 -- novembre 2026
% ============================================================================
% Sources : NIST SP 800-115 (2008) ; PTES ; CIS Controls v8 (contrôle 18).
% ============================================================================

\chapter{Le Hacking Éthique : Définitions, Métiers, Types de Tests}
\label{chap:ethique}
\chaptermark{Le hacking éthique}

\epigraph{%
    « Le testeur et l'attaquant connaissent les mêmes techniques~;\\
    seul le premier rend compte de ce qu'il a fait. »%
}{Principe du cours}

\niveauChapitre{Fondamental}{Prologue~; cours SEC3266.}

\begin{boxobjectifs}
\begin{itemize}
    \item Définir le hacking éthique et le test d'intrusion.
    \item Distinguer test d'intrusion, analyse de vulnérabilités, audit,
          équipe rouge et programme de primes.
    \item Choisir une approche (boîte noire, grise, blanche) et un point de
          départ (externe, interne, présumé compromis).
    \item Décrire les métiers de la sécurité offensive et leurs
          compétences.
    \item Situer la valeur et les limites d'un test.
\end{itemize}
\end{boxobjectifs}

\bigskip

% ============================================================================
\section*{Question directrice : de quoi LiZbiZ avait-elle besoin~: un balayage, un audit ou un test d'intrusion~?}
% ============================================================================

En septembre, M.~KOUNOU propose d'acheter un scanner de vulnérabilités~;
Mme~NKENG rappelle que l'audit d'octobre est prévu~; Mme~MOUAFO veut savoir
si un attaquant peut encore devenir administrateur du domaine. Trois
besoins, trois réponses différentes. Ce chapitre apprend à choisir.

\medskip
\textbf{Plan du chapitre.}
\begin{enumerate}
    \item Définir (section~\ref{sec:c01-definir}).
    \item Les types d'évaluation offensive (section~\ref{sec:c01-types}).
    \item Approches et points de départ (section~\ref{sec:c01-approches}).
    \item Les métiers (section~\ref{sec:c01-metiers}).
    \item Valeur et limites d'un test (section~\ref{sec:c01-limites}).
\end{enumerate}

\separateur

% ============================================================================
\section{Définir}
\label{sec:c01-definir}
% ============================================================================

\begin{boxdef}
\termecle{Hacking éthique}~: emploi des techniques des attaquants, avec
l'autorisation écrite du propriétaire du système, dans un périmètre et une
durée fixés, pour en révéler les faiblesses et aider à les corriger.
\termecle{Test d'intrusion}~: évaluation dans laquelle le testeur imite un
attaquant réel pour contourner les mesures de sécurité et démontrer
l'impact d'une compromission~\cite{nist800115}.
\end{boxdef}

Trois conditions font le caractère éthique~: l'\emph{autorisation} (écrite,
par qui a le pouvoir de la donner), la \emph{limitation} (périmètre,
méthodes, fenêtre) et la \emph{restitution} (tout ce qui est trouvé est
rapporté au client, et à lui seul). Sans l'une d'elles, le même geste devient
une infraction ou une faute professionnelle.

% ============================================================================
\section{Les Types d'Évaluation Offensive}
\label{sec:c01-types}
% ============================================================================

\begin{tableaucours}{L{3.2cm}L{4.6cm}L{4.4cm}}{%
    \tableauHeader{Type} &
    \tableauHeader{Question} &
    \tableauHeader{Livrable}
}
Analyse de vulnérabilités & Quelles failles connues sont présentes~? &
Liste priorisée, souvent automatique \\
Test d'intrusion & Jusqu'où un attaquant peut-il aller~? & Chemins
d'attaque démontrés, constats, recommandations \\
Audit de sécurité & Les mesures sont-elles conformes à un référentiel~? &
Constats d'écarts (cours INF3226) \\
Exercice d'équipe rouge & Détecte-t-on et répond-on à un adversaire
réaliste~? & Scénario, détections, délais \\
Programme de primes & Des chercheurs externes trouvent-ils des failles~? &
Rapports de chercheurs, rémunérés selon la gravité \\
\end{tableaucours}
\vspace{-0.8em}
\begin{center}{\small\itshape Tableau~1.1. Les évaluations offensives}\end{center}

Les CIS Controls consacrent leur contrôle~18 aux tests d'intrusion~: ils
s'adressent aux organisations qui ont déjà mis en place les mesures de
base~; tester avant d'avoir inventorié et corrigé n'apprend rien que le
balayage n'aurait dit~\cite{cis2021controls}.

% ============================================================================
\section{Approches et Points de Départ}
\label{sec:c01-approches}
% ============================================================================

\begin{tableaucours}{L{3.0cm}L{4.8cm}L{4.4cm}}{%
    \tableauHeader{Approche} &
    \tableauHeader{Ce que sait le testeur} &
    \tableauHeader{Intérêt}
}
Boîte noire & Rien, sauf le périmètre & Imite l'attaquant externe~;
coûteuse en temps \\
Boîte grise & Un compte, des plans partiels & Bon compromis, le plus
fréquent \\
Boîte blanche & Code, configurations, schémas & Couverture maximale \\
\end{tableaucours}
\vspace{-0.8em}
\begin{center}{\small\itshape Tableau~1.2. Les approches}\end{center}

Le \termecle{point de départ} compte autant~: \emph{externe} (depuis
Internet), \emph{interne} (depuis un poste du réseau) ou \emph{présumé
compromis} (on suppose qu'un poste est déjà aux mains de l'attaquant, comme
après l'hameçonnage de février). Ce dernier scénario, dit d'\emph{intrusion
présumée}, mesure la défense en profondeur.

\begin{boxscenario}{Opération LOGONE. Le choix de l'approche}
\textbf{Décision} (septembre). Boîte grise~: un compte d'utilisateur
ordinaire et un poste fourni pour l'interne~; les schémas réseau sont
communiqués. Deux points de départ~: externe, et intrusion présumée.

\medskip
\textbf{Analyse} (confiance élevée). La boîte noire aurait consacré des jours
à retrouver ce que LiZbiZ sait déjà~; l'intrusion présumée répond à la vraie
question~: que se passe-t-il quand le prochain hameçonnage réussit~?
\end{boxscenario}

% ============================================================================
\section{Les Métiers}
\label{sec:c01-metiers}
% ============================================================================

\begin{tableaucours}{L{3.4cm}L{8.8cm}}{%
    \tableauHeader{Métier} &
    \tableauHeader{Activité}
}
Pentesteur & Conduit les tests, démontre les chemins d'attaque, rédige les
constats \\
Auditeur de sécurité & Vérifie la conformité aux référentiels, peut
intégrer des tests techniques \\
Analyste SOC et équipe rouge & L'un détecte et répond, l'autre émule
l'adversaire~; ensemble, ils forment l'équipe violette \\
RSSI & Commande, cadre, arbitre les corrections, rend compte à la direction \\
\end{tableaucours}
\vspace{-0.8em}
\begin{center}{\small\itshape Tableau~1.3. Les métiers visés par le cours}\end{center}

\begin{boxmetier}{Pentesteur}
\textbf{Trois compétences au même niveau.} La technique (systèmes, réseaux,
applications), la méthode (périmètre, traçabilité, arrêt) et l'écrit (un
constat qu'un directeur comprend et qu'un administrateur sait corriger).
\end{boxmetier}

% ============================================================================
\section{Valeur et Limites d'un Test}
\label{sec:c01-limites}
% ============================================================================

\begin{boxalert}
\textbf{Violation typique.} Présenter un test sans constat critique comme
une preuve de sécurité.

\smallskip
\textbf{Règle~:} un test est borné dans le temps, le périmètre et les
méthodes~; il prouve l'existence des failles trouvées, pas l'absence des
autres. Le rapport dit ce qui n'a pas été testé.
\end{boxalert}

\begin{boxobjet}{T-00 · Modèle de fiche de constat}
\textbf{Description~:} ce qui a été trouvé, en une phrase. \textbf{Preuve~:}
la trace (capture, requête, horodatage), sans données sensibles.
\textbf{Gravité~:} score CVSS et contexte. \textbf{Impact~:} ce qu'un
attaquant pourrait faire. \textbf{Recommandation~:} la correction et sa
vérification. \textbf{Détection~:} ce qui aurait dû alerter.
\end{boxobjet}

\separateur

\begin{boxresume}
\textbf{Ce qu'il faut retenir de ce chapitre~:}
\begin{itemize}
    \item \textbf{Hacking éthique}~: autorisation, limitation, restitution.
    \item \textbf{Cinq évaluations offensives} répondent à des questions
          différentes.
    \item \textbf{Boîte noire, grise, blanche}~; départ externe, interne ou
          présumé compromis.
    \item \textbf{Quatre métiers}~; l'équipe violette réunit attaque et
          défense.
    \item \textbf{Un test} prouve l'existence d'une faille, pas l'absence des
          autres.
\end{itemize}
\end{boxresume}

\begin{boxinfo}
\textbf{Points de contrôle de l'élève.} Je sais~: définir le hacking éthique
et le test d'intrusion~; choisir un type d'évaluation pour une question
donnée~; choisir une approche et un point de départ~; décrire les métiers~;
énoncer les limites d'un test.
\end{boxinfo}

\bigskip

\begin{boxexo}[title={\ding{45}\quad Exercice~1.1. Quel type~? \niveaubase}]
Quel type d'évaluation répond à~: « nos serveurs ont-ils des correctifs en
retard~? »~; « notre SOC verrait-il un attaquant~? »~; « sommes-nous
conformes à l'ISO/IEC~27001~? »
\end{boxexo}

\begin{boxexo}[title={\ding{45}\quad Exercice~1.2. Les trois conditions \niveaubase}]
Un étudiant teste le site de son école « pour l'aider ». Quelles conditions
du hacking éthique manquent~?
\end{boxexo}

\begin{boxexo}[title={\ding{45}\quad Exercice~1.3. Approche \niveaubase}]
Quelle approche choisir pour l'ANPP, qui fournit le code de son portail~?
\end{boxexo}

\begin{boxexo}[title={\ding{45}\quad Exercice~1.4. Intrusion présumée \niveaumoyen}]
Expliquez en quoi le point de départ « présumé compromis » répond mieux à la
question de LiZbiZ qu'un test externe seul.
\end{boxexo}

\begin{boxexo}[title={\ding{45}\quad Exercice~1.5. Contrôle 18 \niveaumoyen}]
Pourquoi les CIS Controls placent-ils les tests d'intrusion en dernier~?
Que doit faire une organisation avant~?
\end{boxexo}

\begin{boxexo}[title={\ding{45}\quad Exercice~1.6. Violette \niveaumoyen}]
Qu'apporte une équipe violette par rapport à un test d'intrusion classique~?
\end{boxexo}

\begin{boxexo}[title={\ding{45}\quad Exercice~1.7. Primes \niveaumoyen}]
Quels prérequis une organisation doit-elle réunir avant d'ouvrir un
programme de primes aux failles~?
\end{boxexo}

\begin{boxexo}[title={\ding{45}\quad Exercice~1.8. Limites \niveaumoyen}]
Rédigez le paragraphe « limites » du rapport de l'Opération LOGONE.
\end{boxexo}

\begin{boxexo}[title={\ding{45}\quad Exercice~1.9. Un programme \niveauexpert}]
Proposez à la BANQUE SANAGA un programme d'évaluations offensives sur deux
ans, en combinant les types du Tableau~1.1.
\end{boxexo}

\begin{boxexo}[title={\ding{45}\quad Exercice~1.10. Le profil \niveauexpert}]
Rédigez la fiche de poste d'un pentesteur junior pour NYONG SÉCURITÉ~:
missions, compétences, éthique, progression.
\end{boxexo}

\bigskip

\begin{boxinfo}
\textbf{Pour aller plus loin.}
\begin{itemize}
    \item NIST SP~800-115~\cite{nist800115}~; CIS Controls~\cite{cis2021controls}.
\end{itemize}
\end{boxinfo}

\begin{boxinfo}
\textbf{Transition.} L'autorisation n'est pas une formalité~: c'est elle
que la loi regarde. Le Chapitre~\ref{chap:droit} présente le droit
camerounais.
\end{boxinfo}
